% Contextual printed508: blank reverse; principal515, Google519.
\par\clearpage\null\clearpage
% Unnumbered printed509; principal516, Google520.
\begin{center}
{\Large\textbf{\textls[160]{POSTHUMOUS PAPERS.}}}\par\bigskip
{\large [Reduction of oblique angles to the horizon.]}\par\bigskip
{[}1.]\par\medskip
{\large Reduction to the horizon.}
\end{center}
\[
\begin{array}{ll}
{[}A\ .\ .\ .\ .\ .\ .\ .&\text{oblique angle}\\
A+x\ .\ .\ .\ .&\text{horizontal angle}\\
h,\,h'\ .\ .\ .\ .\ .&\text{elevations.}
\end{array}
\]
\noindent Set]
\[
\begin{gathered}
\tfrac12(A+h+h')=s\\
\tfrac12(h+h')=\sigma,\qquad\tfrac12(h-h')=\delta;
\end{gathered}
\]
\noindent [then]
\[
\begin{aligned}
\operatorname{tang}\tfrac12(A+x)
&=\sqrt{\frac{\sin(\frac12A+\delta)\sin(\frac12A-\delta)}{\cos(\frac12A+\sigma)\cos(\frac12A-\sigma)}}\\
&=\sqrt{\frac{\sin(s-h)\sin(s-h')}{\cos s\cos(A-s)}}.
\end{aligned}
\]
\begin{center}[2.]\par\medskip
{\large Reduction of oblique angles to the horizon.}\end{center}
\[
\begin{array}{ll}
A\ .\ .\ .\ .\ .\ .&\text{oblique angle}\\
x\ .\ .\ .\ .\ .\ .&\text{correction}\\
h,\,h'\ .\ .\ .\ .&\text{elevations}
\end{array}
\]
\[
\frac{\operatorname{tang}\frac12A}{4\,.\,206265}(h+h')^2=M,\qquad
\frac{\operatorname{cotang}\frac12A}{4\,.\,206265}(h-h')^2=N.
\]
\noindent Approximate value $x=M-N$.

% Printed510; principal517, Google521.
\par\clearpage
\noindent\makebox[0pt][l]{\fontsize{13}{16}\selectfont 510}\hfill{\small\textbf{POSTHUMOUS PAPERS.}}\hfill\null\par\medskip
\[
\begin{aligned}
\sec\tfrac12x^{\frac13}\,.\,\frac{\sin(A+x)}{\sin(A+\frac12x)}&=\alpha\\[4pt]
\frac{\operatorname{tang}\frac12(A+x)}{\operatorname{tang}\frac12A}&=\beta\\[4pt]
\left(\sec\frac{h+h'}2\right)^{\frac23}&=\mu\\[4pt]
\left(\sec\frac{h-h'}2\right)^{\frac23}&=\nu.
\end{aligned}
\]
\noindent The corrected value is $\displaystyle x=\frac{\alpha\beta M}{\mu}-\frac{\alpha N}{\beta\nu}$.
\begin{center}Example.\end{center}
\noindent Oblique angle measured between Clausberg and Geismar
\[
A=105^\circ27'0\overset{\prime\prime}{,}78,\qquad
\tfrac12A=52^\circ43'30\overset{\prime\prime}{,}39;
\]
\[
\begin{array}{r@{\;}c@{\;}r@{\;}r@{\;}r@{\;}c@{\;}r}
h&=&2^\circ&19'&39\overset{\prime\prime}{,}75&=&8379\overset{\prime\prime}{,}75\\
h'&=&1&4&3{,}17&=&3843{,}17.
\end{array}
\]
\noindent\begin{minipage}[t]{.31\linewidth}\centering
$\begin{array}{r@{\ }r}
h+h'\ .\ .\ .\ .\ .&\underline{4{,}087175}\\
{[}(h+h')^2\ .\ .\ .&8{,}174350\\
\frac{\operatorname{tang}\frac12A}{4\,.\,206265}]\ .\ .\ .&\underline{4{,}202072}\\
M\ .\ .\ .\ .\ .\ .\ .&2{,}376422\\[3pt]
\beta\ .\ .\ .\ .\ .\ .\ .\ .&+477\\
1:\mu\ .\ .\ .\ .\ .\ .&-127\\
\alpha\ .\ .\ .\ .\ .\ .\ .\ .&\underline{-\phantom{1}64}\\
&2{,}376708\\[3pt]
&238\overset{\prime\prime}{,}07
\end{array}$
\end{minipage}\hfill
\begin{minipage}[t]{.31\linewidth}\centering
$\begin{array}{r@{\ }r}
h-h'\ .\ .\ .\ .\ .&\underline{3{,}656728}\\
{[}(h-h')^2\ .\ .\ .&7{,}313456\\
\frac{\operatorname{cotang}\frac12A}{4\,.\,206265}]\ .\ .\ .&\underline{3{,}964958}\\
N\ .\ .\ .\ .\ .\ .\ .&1{,}278414\\[3pt]
1:\beta\ .\ .\ .\ .\ .\ .&-477\\
\makebox[0pt][r]{$-$\qquad}1:\nu\ .\ .\ .\ .\ .\ .&-\phantom{1}17\\
\alpha\ .\ .\ .\ .\ .\ .\ .\ .&\underline{-\phantom{1}64}\\
&1{,}277856\\[3pt]
&18\overset{\prime\prime}{,}96
\end{array}$
\end{minipage}\hfill
\begin{minipage}[t]{.34\linewidth}\centering
$\begin{array}{r@{\ }r}
{[}1:4\,.\,206265]\ .\ .\ .&4{,}083515\\
\operatorname{tang}\tfrac12A\ .\ .\ .\ .\ .\ .\ .&0{,}118557\\
\operatorname{tang}\tfrac12(A+x)\ .\ .\ .&\underline{0{,}119034}\\
{[}\sqrt[3]{\sec\tfrac12x}\ .\ .\ .\ .\ .\ .\ .&0{,}000000\\
\sin(A+x)\ .\ .\ .\ .\ .&9{,}983887\\
\sin(A+\tfrac12x)\ .\ .&\underline{9{,}983951]}\\
M=&237\overset{\prime\prime}{,}91\\
N=&\underline{18{,}99}\\
M-N=&218{,}92
\end{array}$
\end{minipage}
\[
x=+219\overset{\prime\prime}{,}11.\hspace*{.36\linewidth}
\]
\noindent Horizontal angle $=105^\circ30'39\overset{\prime\prime}{,}89$.

% Printed511; principal518, Google522.
\par\clearpage
\noindent\null\hfill{\small\textbf{REDUCTION OF OBLIQUE ANGLES TO THE HORIZON.}}\hfill\makebox[0pt][r]{\fontsize{13}{16}\selectfont 511}\par\bigskip
\begin{center}[3].\par\medskip
{\large Reduction of an oblique angle to the horizon.}\end{center}
\noindent [Let]
\[
\begin{array}{ll}
\alpha\ \text{and}\ \beta\ .\ .\ .\ .\ .\ .&\text{elevations expressed in arcseconds}\\
A\ .\ .\ .\ .\ .\ .\ .\ .\ .\ .\ .\ .&\text{[measured] angular distance}\\
A'=A+r\ .\ .\ .\ .&\text{horizontal angular distance.}
\end{array}
\]
\noindent [Then]
\[
2\sin\tfrac12r=\frac{\sin\frac12(\alpha+\beta)^2\operatorname{tang}\frac12A-\sin\frac12(\alpha-\beta)^2\operatorname{cotang}\frac12A}{\cos\alpha\cos\beta}\,.\,\frac{\sin A}{\sin(A+\frac12r)}.
\]
\noindent [Since the approximation
\[
\sin\tfrac12u=\frac{u''}{2\,.\,206265}\sqrt[3]{\cos\tfrac12u}
\]
\noindent holds, the following scheme can be used to calculate $r$.]
\[
\begin{array}{r@{\ }l}
\mathrm C\log\cos\alpha=&\\
\mathrm C\log\cos\beta=&\\
\log\mathrm{const}=[\log\frac1{4\,.\,206265}=&\underline{4{,}083\,5149-10]}\\
{[}\log c=]&
\end{array}
\]
\noindent\begin{minipage}[t]{.49\linewidth}\centering
$\begin{array}{r@{\ }l}
\log\operatorname{tang}\tfrac12A=&\\
2\log(\alpha+\beta)=&\\
\tfrac23\log\cos\tfrac12(\alpha+\beta)=&\underline{\hspace{2.7cm}}\\
\log a=&\\
a=&
\end{array}$
\end{minipage}\hfill
\begin{minipage}[t]{.49\linewidth}\centering
$\begin{array}{r@{\ }l}
\log\operatorname{cotang}\tfrac12A=&\\
2\log(\alpha-\beta)=&\\
\tfrac23\log\cos\tfrac12(\alpha-\beta)=&\underline{\hspace{2.2cm}}\\
\log b=&\\
b=&
\end{array}$
\end{minipage}
\par\bigskip
\noindent\begin{minipage}[t]{.34\linewidth}\centering $a-b=$\end{minipage}\hfill
\begin{minipage}[t]{.65\linewidth}\raggedleft
$\begin{array}{r@{\ }l}
\log(a-b)=&\\
{[}\log c=]&\\
\mathrm{Corr.}=\left[\log\left\{\sec\tfrac12r^{\frac13}\,.\,\frac{\sin A}{\sin(A+\frac12r)}\right\}=\right]&\underline{\hspace{1.8cm}}\\
\log r=&\\
r=&
\end{array}$
\end{minipage}
\begin{center}\vspace{1cm}\rule{2.7cm}{.4pt}\end{center}

% Printed512; principal519, Google523.
\par\clearpage
\noindent\makebox[0pt][l]{\fontsize{13}{16}\selectfont 512}\hfill{\small\textbf{REMARKS.}}\hfill\null\par\bigskip
\begingroup\small
\begin{center}\vspace{2cm}\textbf{REMARKS.}\end{center}
\par\noindent\hspace*{\parindent}Note [1] is found in a journal of the observatory, II 1813--1817, note [2] on the last page of a table of logarithms, and note [3] on the last page of \textsc{Gauss}'s copy of »G. F. \textsc{Röslers} Handbuch der praktischen Astronomie für Anfänger und Liebhaber, Erster Theil. Tübingen 1788.«
\par\noindent\hspace*{\parindent}In [2] the original gives the expression for $\alpha$
\[
\alpha=\sec x^{\frac13}\,.\,\frac{\sin(A+x)}{\sin(A+\frac12x)},
\]
\noindent which should read as given above. A few small errors in the numerical example have been corrected.
\par\noindent\hspace*{\parindent}The formula underlying note [2] can be derived as follows. Let $A$ be the oblique angle, $A+x$ the corresponding horizontal angle, and $h$ and $h'$ the elevation angles, so that $90^\circ-h$ and $90^\circ-h'$ are the zenith distances. Then
\[
\begin{aligned}
\text{1)}\quad\cos A&=\sin h\sin h'+\cos h\cos h'\cos(A+x)\\
&=\sin h\sin h'\{\cos\tfrac12(A+x)^2+\sin\tfrac12(A+x)^2\}\\
&\qquad+\cos h\cos h'\{\cos\tfrac12(A+x)^2-\sin\tfrac12(A+x)^2\}\\
&=\cos\tfrac12(A+x)^2\cos(h-h')-\sin\tfrac12(A+x)^2\cos(h+h');
\end{aligned}
\]
\noindent and subtracting from this
\[
\cos(A+x)=\cos\tfrac12(A+x)^2-\sin\tfrac12(A+x)^2,
\]
\noindent gives the equation
\[
\sin\tfrac12x\sin(A+\tfrac12x)=-\cos\tfrac12(A+x)^2\sin\tfrac12(h-h')^2+\sin\tfrac12(A+x)^2\sin\tfrac12(h+h')^2
\]
\noindent or
\[
\text{2)}\quad\sin\tfrac12x\,.\,\frac{\sin(A+\frac12x)}{\sin(A+x)}
=\tfrac12\operatorname{tang}\tfrac12(A+x)\sin\tfrac12(h+h')^2-\tfrac12\operatorname{cotang}\tfrac12(A+x)\sin\tfrac12(h-h')^2.
\]
\par\noindent\hspace*{\parindent}Thus, to a first approximation, if $x$, $h$ and $h'$ are expressed in arcseconds and we set $\rho=\frac1{206265}$, we obtain
\[
\begin{aligned}
x&=\tfrac14\rho(h+h')^2\operatorname{tang}\tfrac12A-\tfrac14\rho(h-h')^2\operatorname{cotang}\tfrac12A\\
\text{3)}\qquad&=M-N.
\end{aligned}
\]
\par\noindent\hspace*{\parindent}Since, for small angles, the approximation
\[
\sin u=\rho\,.\,u''\cos u^{\frac13}
\]
\par\endgroup

% Printed513; principal520, Google524.
\par\clearpage
\noindent\null\hfill{\small\textbf{REDUCTION OF OBLIQUE ANGLES TO THE HORIZON.}}\hfill\makebox[0pt][r]{\fontsize{13}{16}\selectfont 513}\par\medskip
\begingroup\small
\noindent holds, equation (2) gives the more accurate approximation
\[
\text{4)}\quad x=\left(\sec\frac x2\right)^{\frac13}\,.\,\frac{\sin(A+x)}{\sin(A+\frac12x)}
\GaussOldenburgBracePair{\tfrac14\rho\frac{(h+h')^2}{\left(\sec\frac{h+h'}2\right)^{\frac23}}\operatorname{tang}\tfrac12(A+x)
-\tfrac14\rho\frac{(h-h')^2}{\left(\sec\frac{h-h'}2\right)^{\frac23}}\operatorname{cotang}\tfrac12(A+x)}
\]
\noindent or, with the notation at the top of p. 510,
\[
x=\alpha\left\{\frac{\beta M}{\mu}-\frac N{\beta\nu}\right\}.
\]
\par\bigskip\noindent\hspace*{\parindent}The formula in note [3] is a little more convenient to use, since it requires less indirect calculation. It can be derived as follows. By equation (1),
\[
\cos(A+x)=\frac{\cos A-\sin h\sin h'}{\cos h\cos h'};
\]
\noindent hence
\[
\begin{aligned}
\cos(A+x)-\cos A=\frac1{\cos h\cos h'}\bigl\{&(\cos\tfrac12A^2-\sin\tfrac12A^2)(1-\cos h\cos h')\\
&-(\cos\tfrac12A^2+\sin\tfrac12A^2)\sin h\sin h'\bigr\}
\end{aligned}
\]
\noindent or
\[
-\sin(A+\tfrac12x)\sin\tfrac12x=\frac1{\cos h\cos h'}\{\cos\tfrac12A^2\sin\tfrac12(h-h')^2-\sin\tfrac12A^2\sin\tfrac12(h+h')^2\}
\]
\noindent or
\[
\text{5)}\quad\sin\tfrac12x\frac{\sin(A+\frac12x)}{\sin A}
=\frac1{2\cos h\cos h'}\{\operatorname{tang}\tfrac12A\sin\tfrac12(h+h')^2-\operatorname{cotang}\tfrac12A\sin\tfrac12(h-h')^2\}.
\]
\par\noindent\hspace*{\parindent}Putting $\sin u=\rho\,.\,u''\cos u^{\frac13}$ again gives
\[
\text{6)}\quad x=\frac{\rho\left(\sec\frac x2\right)^{\frac13}}{4\cos h\cos h'}\,.\,\frac{\sin A}{\sin(A+\frac12x)}
\GaussOldenburgBracePair{\frac{(h+h')^2}{\left(\sec\frac{h+h'}2\right)^{\frac23}}\operatorname{tang}\tfrac12A
-\frac{(h-h')^2}{\left(\sec\frac{h-h'}2\right)^{\frac23}}\operatorname{cotang}\tfrac12A}
\]
\noindent or, if we also put
\[
\text{7)}\qquad\alpha'=\left(\sec\frac x2\right)^{\frac13}\,.\,\frac{\sin A}{\sin(A+\frac12x)}
\]
\noindent and, as before,
\[
\begin{gathered}
M=\rho\left(\frac{h+h'}2\right)^2\operatorname{tang}\tfrac12A,\qquad N=\rho\left(\frac{h-h'}2\right)^2\operatorname{cotang}\tfrac12A\\
\mu=]\sqrt[3]{\left(\sec\frac{h+h'}2\right)^2},\qquad\nu=\sqrt[3]{\left(\sec\frac{h-h'}2\right)^2}
\end{gathered}
\]
\noindent we obtain
\[
\text{8)}\qquad x=\frac{\alpha'}{\cos h\cos h'}\left\{\frac M\mu-\frac N\nu\right\}.
\]
\begin{flushright}\textsc{Krüger.}\end{flushright}
\par\endgroup
\begin{center}\rule{3.5cm}{.4pt}\end{center}
\noindent\makebox[0pt][l]{\small IX.}\makebox[.86\linewidth][r]{65}\par

% Unnumbered printed514; principal521, Google525.
\par\clearpage
\begin{center}
{\Large POSTHUMOUS PAPERS AND CORRESPONDENCE.}\par\vspace{1cm}
{\large [Observations by repetition.]}\par\bigskip
{[}1.]
\end{center}
\[
\frac{\text{Probable error of reading}}{\text{Probable error of pointing}}=\frac1k.
\]
\noindent [Let $A$, $B$, $C$,\ .\ .\ .\ be the equally spaced readings.]
\begin{center}Most probable result\end{center}
\noindent from observations
\[
\begin{array}{c@{\qquad\qquad}l}
3&\dfrac{(kk+3)(D-A)+(C-B)}{3kk+10}\\[10pt]
4&\dfrac{(kk+2)(E-A)+(D-B)}{4kk+10}\\[10pt]
5&\dfrac{(k^4+5kk+5)(F-A)+(kk+3)(E-B)+(D-C)}{5k^4+28kk+35}\\[10pt]
6&\dfrac{(k^4+4kk+3)(G-A)+(kk+2)(F-B)+(E-C)}{6k^4+28kk+28}.
\end{array}
\]
\par\noindent\hspace*{\parindent}The coefficients form a recurring series with scale $kk+2$, $-1$.
\par\bigskip\noindent\hspace*{\parindent}If readings are taken only at the beginning and after the first, $(n-1)^{\text{th}}$, and $n^{\text{th}}$ observations [and these readings are $A$, $B$, $Y$, $Z$], then the most probable result
\[
=\frac{((n-2)kk+n)(Z-A)+(n-2)(Y-B)}{n(n-2)kk+nn+(n-2)^2}.
\]
\par\noindent\hspace*{\parindent}Probable error
\[
=\sqrt{\frac{(n-2)k^4+(2n-2)kk+2}{(nn-2n)k^4+(2nn-4n+4)kk}}.
\]

% Printed515; principal522, Google526.
\par\clearpage
\noindent\null\hfill{\fontsize{8.2}{10}\selectfont\textbf{POSTHUMOUS PAPERS AND CORRESPONDENCE.\quad OBSERVATIONS BY REPETITION.}}\hfill\makebox[0pt][r]{\fontsize{13}{16}\selectfont 515}\par\medskip
\par\noindent\hspace*{\parindent}The probable .error when readings are taken \textls[120]{only} at the beginning and end is
\[
\sqrt{\frac{nkk+2}{nnkk}};
\]
\noindent by \textsc{Svanberg}'s method,
\[
\frac{60+6(n(n+2)+2)kk}{5n(n+1)(n+2)kk}.
\]
\begin{center}\vspace{.7cm}[2.]\par\bigskip
{\large\textsc{Gauss} to \textsc{Bessel}.\quad Göttingen, 27 January 1819.}\end{center}
\par\noindent\hspace*{\parindent}.\ .\ .\ .\ .\ .\enspace I am now considering some further additions to my theory of least squares. One point concerns the treatment of observations made with repeating instruments. \textsc{Laplace} has recently published a paper on this and, among other things, criticized \textsc{Svanberg}'s procedure. Yet what \textsc{Laplace} says of \textsc{Svanberg}, »c’est un nouvel exemple des illusions auxquelles on est exposé dans ces recherches délicates«, might in a certain sense be applied to that great geometer himself. The procedure of taking the whole arc traversed and simply dividing it by $n$ (which \textsc{Laplace} defends) agrees with the theory of probability only if the reading (and graduation) errors are negligible in comparison with the pointing errors, especially if the number of observations is not very large. \textsc{Laplace} completely ignored \textls[120]{the former} errors. It is very remarkable that if, conversely, the pointing errors are negligible in comparison with the reading and graduation errors, the procedure in accordance with the theory of probability is exactly the same as \textsc{Svanberg}'s. It seems that \textsc{Svanberg} had only the latter errors in mind (which in French instruments may well considerably exceed the pointing errors). If both kinds of error are considered to coexist, the problem seems very difficult, but approached in the right way it is extremely simple and elementary, as you will no doubt see at once. Yet, like Columbus's egg, this way of proceeding may well fail to occur to everyone immediately. I observe more generally that some astronomers, for example \textsc{Lindenau} and \textsc{Littrow}, apply the method
\par\medskip\noindent\makebox[.86\linewidth][r]{65\textsuperscript{*}}\par

% Printed516; principal523, Google527.
\par\clearpage
\noindent\makebox[0pt][l]{\fontsize{13}{16}\selectfont 516}\hfill{\small\textbf{CORRESPONDENCE.\quad OBSERVATIONS BY REPETITION.}}\hfill\null\par\medskip
\noindent of least squares without always doing so in its true spirit. For this reason it seems to me that there may be some merit in developing several points of this theory separately. For 12 observations, when the graduation errors are 3 times as large as the pointing errors, the weight of the result is
\[
\begin{array}{l@{\;}c@{\;}l}
\text{by \textsc{Laplace}'s method}&=&4{,}80\\
\text{by \textsc{Svanberg}'s method}&=&7{,}00\\
\text{by the rigorous method}&=&7{,}08,
\end{array}
\]
\noindent with the weight from one observation \textls[120]{without} reading error set $=1$.
\par\noindent\hspace*{\parindent}Merely including the first and eleventh readings would raise the weight under the correct method to 6,26.\ .\ .\ .\ .\ .\ .
\begin{center}\rule{2.7cm}{.4pt}\end{center}
\begingroup\small
\begin{center}\textbf{REMARKS.}\end{center}
\par\noindent\hspace*{\parindent}The note on observations by repetition is written on the last leaf of \textsc{Gauss}'s copy of \textsc{Laplace}'s »Deuxième supplément à la théorie analytique des probabilités, Février 1818«. The letter to \textsc{Bessel} excerpted under [2], printed from the original, also refers to this paper. A theory of observations by repetition, for both equally and unequally spaced readings, was published in 1834 by \textsc{Bessel} (Astronomische Nachrichten, vol. XI, no. 256, pp. 269 ff.).
\par\noindent\hspace*{\parindent}Let the weight of a reading error be 1 and that of a single sighting error be $p$. If readings are taken at regular intervals of $r$ repetitions, then the error equal to the sum of the sighting errors between 2 successive readings has weight $\frac{p}{2r}$. \textsc{Gauss} puts $\dfrac{p}{2r}=\dfrac1{kk}$.
\par\noindent\hspace*{\parindent}Let $x$ denote the angle sought, so that $rx$ is obtained from each set of 2 successive readings. Let the readings themselves be $l_0$, $l_1$,\ldots, $l_n$. Further, let $\mu_1$, $\mu_2$, $\mu_3$,\ldots\ be a sequence of numbers obtained from the recurrence formula
\[
-\mu_{i-1}+g\mu_i-\mu_{i+1}=0,\qquad g=\frac{2r}p+2=kk+2
\]
\noindent subject to the conditions $\mu_0=0$ and $\mu_1=1$. Thus
\[
\begin{array}{r@{\;}c@{\;}l@{\;}c@{\;}l}
\mu_1&=&1&&\\
\mu_2&=&g&=&kk+2\\
\mu_3&=&gg-1&=&k^4+4kk+3\\
\mu_4&=&g^3-2g&=&k^6+6k^4+10kk+4\\
\multicolumn{5}{c}{\text{etc.}}
\end{array}
\]
\par\noindent\hspace*{\parindent}The $\mu$ can be computed directly from the formula
\[
k\sqrt{kk+4}\,.\,\mu_i=u^i-u^{-i};\qquad u=\tfrac12(kk+2+k\sqrt{kk+4}).
\]
\par\endgroup

% Printed517; principal524, Google528.
\par\clearpage
\noindent\null\hfill{\small\textbf{REMARKS.\quad OBSERVATIONS BY REPETITION.}}\hfill\makebox[0pt][r]{\fontsize{13}{16}\selectfont 517}\par\medskip
\begingroup\small
\par\noindent\hspace*{\parindent}If the number of readings is $n=2m$, then in general (according to \textsc{Andrae} and \textsc{Zachariae}, Den danske Gradmaaling, vol. II, 1872, pp. 296 ff.)
\[
x=\frac2p\,.\,\frac{\displaystyle\sum_{\lambda=0}^{m-1}\mu_{m-\lambda}(l_{n-\lambda}-l_\lambda)}{n\mu_{m+1}-(n+2)\mu_m}
\]
\noindent or equivalently
\[
rx=\frac{\displaystyle\sum_{\lambda=0}^{m-1}\mu_{m-\lambda}(l_{n-\lambda}-l_\lambda)}{\dfrac1{kk}\{(nkk+n-2)\mu_m-n\mu_{m-1}\}}.
\]
\par\noindent\hspace*{\parindent}The reciprocal weight of $x$ is
\[
\frac1P=\frac2{pr}\,.\,\frac{\mu_{m+1}-\mu_m}{n\mu_{m+1}-(n+2)\mu_m}
=\frac{kk}{rr}\left(\frac1n+\frac1n\,.\,\frac{2\mu_{\frac n2}}{(nkk+n-2)\mu_{\frac n2}-n\mu_{\frac n2-1}}\right).
\]
\par\noindent\hspace*{\parindent}Thus if we set
\[
(nkk+n-2)\mu_m-n\mu_{m-1}=kkM,
\]
\noindent then
\[
\frac1P=\frac{kk}{rr}\left(\frac1n+\frac{2\mu_m}{nkkM}\right)
=\frac1{rrM}((kk+1)\mu_m-\mu_{m-1})=\frac1{rrM}(\mu_{m+1}-\mu_m).
\]
\par\noindent\hspace*{\parindent}And if the number of readings is $n=2m+1$, then
\[
x=\frac2p\,.\,\frac{\displaystyle\sum_{\lambda=0}^{m}(\mu_{m-\lambda+1}+\mu_{m-\lambda})(l_{n-\lambda}-l_\lambda)}{n\mu_{m+2}-2\mu_{m+1}-(n+2)\mu_m}
\]
\noindent or
\[
rx=\frac{\displaystyle\sum_{\lambda=0}^{m}(\mu_{m-\lambda+1}+\mu_{m-\lambda})(l_{n-\lambda}-l_\lambda)}{\dfrac1{kk}\{(nkk+2n-2)\mu_{m+1}-2(n+1)\mu_m\}}.
\]
\par\noindent\hspace*{\parindent}Here the reciprocal weight of $x$ is
\[
\begin{aligned}
\frac1P&=\frac2{pr}\,.\,\frac{\mu_{m+2}-\mu_m}{n\mu_{m+2}-2\mu_{m+1}-(n+2)\mu_m}\\
&=\frac{kk}{rr}\left(\frac1n+\frac1n\,.\,\frac{2(\mu_{\frac{n+1}2}+\mu_{\frac{n-1}2})}{(nkk+2n-2)\mu_{\frac{n+1}2}-2(n+1)\mu_{\frac{n-1}2}}\right).
\end{aligned}
\]
\par\noindent\hspace*{\parindent}Thus if we set
\[
(nkk+2n-2)\mu_{m+1}-2(n+1)\mu_m=kkN
\]
\noindent then
\[
\frac1P=\frac{kk}{rr}\left(\frac1n+\frac{2(\mu_{m+1}+\mu_m)}{nkkN}\right)
=\frac1{rrN}((kk+2)\mu_{m+1}-2\mu_m)=\frac1{rrN}(\mu_{m+2}-\mu_m).
\]
\par\noindent\hspace*{\parindent}Writing $A$, $B$, $C$,\ldots\ in place of $l_0$, $l_1$, $l_2$,\ldots, we obtain for $n=4$ and 6, and for $n=3$ and 5, the results listed on p. 514, which therefore apply to the $r$-fold angle. The $\mu$, and hence the coefficients in the numerator of the expression for the $r$-fold angle when $n$ is even, are the coefficients of the expansion of the fraction $\dfrac1{1-(kk+2)y+yy}$ into a series in ascending powers of $y$. If $n$ is odd, the coefficients in the numerator of the fraction giving the $r$-fold angle are those of the series obtained by expanding $\dfrac{1+y}{1-(kk+2)y+yy}$.
\par\endgroup
